\scanpage{52}
\mspage{8}
% archive sheet number: 51
On a ainsi la notion de \uline{graphe pondéré} \deleted{\unsure{orienté}}\inserted{\uline{à repères}}
\margin{$\unsure{r|\Gamma}$}

\diagram{\begin{tikzcd}[ampersand replacement=\&]
{\vec R_\ell^+\,\deleted{\unsure{\mathbb D_\infty}}}\arrow[r,phantom,"\subset"]\&\underset{\sigma_0,\sigma_1}{\vec R}\arrow[r]\&\vec A\arrow[loop right,"\sigma"]\arrow[d,"o"']\&\\
\&\&S\arrow[r,"\unsure{\gamma|g}"]\&\mathbb N
\end{tikzcd}}
\margin{Groupoïde (\uline{Gr.pondrep})\par foncteur}
\diagram{\begin{tikzcd}[ampersand replacement=\&]
S_*\mathscr T\unsure{\mathbb D|D}_{\mathrm{top}}^{\pm}\arrow[r]\&(\mathrm{Grpond\ rep})
\end{tikzcd}}
qui sera importante sans doute dans l'explicitation de l'opération de ``recollement''. Il est clair en \unsure{tous|tout} cas que toute description-traduction algébrico-combinatoire de $\mathscr T\mathbb D_{\mathrm{top}}^{\pm}$ devra permettre, pour chaque objet du \deleted{la catégorie}\inserted{groupoïde} paradigme, de reconstituer (à isom. unique près) le \uline{graphe pondéré à repères} correspondant. -- plus précisément, il faut reconstituer le foncteur canonique
\diagram{\begin{tikzcd}[ampersand replacement=\&]
S_*\mathscr T\mathbb D_{\mathrm{top}}^{\pm}\arrow[r]\&(\mathrm{Gr.pond.rep.})
\end{tikzcd}}
(et non seulement le foncteur à valeurs dans $(\mathrm{Grbipond})$)

\uline{Opérations essentielles}

A) Opérations ``d'oubli'' ou de ``\uline{gommage}''.

\textcircled{1} \uline{Gommage des repères} À un objet $(X,\Sigma,R)$, muni en plus d'une partie $R'\subset R$, il associe $(X,\Sigma,R\setminus R')$ (gommage de $R'$).
